\section{Build vs. Buy：差异化边界与全球 control plane}

平台公司很容易陷入“既然基础设施重要，就全部自研”的误区。Rauch 的判断相反：AWS 的区域、骨干网、对象存储和物理运营规模难以复制，Vercel 更应该在 application-aware layer 积累差异化。Build vs. Buy 不是技术能力测试，而是资本、时间和学习效应的配置问题。

\subsection{把自研放在 insight 会复利的位置}

上一节已经把 hyperscaler primitives 与 application intent 之间的缺口解释为“编译层问题”；本节进一步回答组织边界：哪些层值得自己掌握，哪些层应该借助成熟供应商。读图时先比较 commodity scale 与 application-aware insight，再判断某项自研是否会持续产生产品知识、控制面能力和单位经济性优势。

\lecturefigure{06-build-vs-buy.png}{复用 commodity scale，在 application-aware layer 建立专用能力。}{本地字幕 07:40--10:30；概念重绘。}

\begin{knowledgebox}{读图：Apple silicon 类比的边界}
课堂用 Apple 最终自研芯片类比 vertical integration：只有当某层成为产品体验和长期经济性的核心，且组织拥有足够规模与能力时，向下整合才可能合理。类比不能直接推出 Vercel 应造数据中心；它只说明 build/buy boundary 会随着战略变化移动。
\end{knowledgebox}

\subsection{全球 metadata 才是平台差异化}

Vercel 更像 application-aware CDN/control plane，而不是 EC2。每个 deployment 是不可变对象；domain binding 是一个可变指针；全球 edge 需要快速获得最新 pointer 和 routing metadata。回滚不需要覆写旧部署，只需把 domain 指针切回已知版本。

\lecturefigure{07-global-metadata.png}{不可变部署与可变 domain pointer 分离了 artifact identity 和流量决策。}{本地字幕 10:00--16:40、33:30--34:05；概念重绘。}

\begin{importantbox}{Control plane 与 data plane}
\term{control plane} 决定 deployment、domain、routing、policy 和 rollout；\term{data plane} 真正处理用户请求、执行函数和返回 assets。控制面更新必须快速传播，但不应让单点故障中断已知路由；数据面必须在 metadata 短暂陈旧时保持安全行为。
\end{importantbox}

\subsection{Opinionated platform 的正负面}

框架约定提供 guard rails，使平台能推断 routing、cache 和 runtime。代价是部分自由被收敛，非典型 workload 可能需要 escape hatch。一个健康的 opinionated platform 应让常见路径极简，同时保留 artifact contract、标准协议和 override，防止生态锁死。

\lecturefigure{08-opinionated-platform.png}{框架约定为平台创造可优化的结构化表面。}{Vercel FDI 官方材料；本地字幕 13:30--17:15。}

\begin{warningbox}{Next.js 不是“把用户锁进 Vercel”的充分证据}
框架与平台之间确实存在正反馈和商业利益，但判断 lock-in 应检查开放源码、标准输出、self-host 路径、数据可迁移性和替代成本。讲义讨论 mechanism，不采用“Trojan horse”式情绪标签替代分析。
\end{warningbox}

\teachervoice{课堂没有把 Vercel 与 AWS 描述成零和竞争。AWS 提供可信的 hyperscale primitives，Vercel 把尚未迁移到云的应用通过更高层模型接入这些 primitives。平台之间既竞争价值份额，也可能互相扩大市场。}

\subsection{本章小结}

Build vs. Buy 应围绕差异化和学习复利；Vercel 的专用能力集中在 framework compiler、global metadata、routing 和 developer loop，而不是重复 hyperscaler 的物理规模。

